\section{Conditional Compilation}
\label{sec:cond_compilation}
\index{directive syntax!conditional compilation}
\index{conditional compilation!_OPENMP macro@\kcode{_OPENMP} macro}
\index{conditional compilation!sentinel}

The OpenMP API supports conditional compilation by defining the 
\kcode{_OPENMP} macro name for preprocessing.
In addition, for Fortran the following conditional compilation sentinels
are recognized:
\begin{boxedcode}
!$               ! in Fortran free form sources
!$ | *$ | c$     ! in Fortran fixed form sources
\end{boxedcode}

\begin{ccppspecific}
The following example illustrates the use of conditional compilation using 
the OpenMP macro \kcode{_OPENMP}. With OpenMP compilation, the \kcode{_OPENMP} 
macro becomes defined.

\cnexample{cond_comp}{1}
\end{ccppspecific}

\begin{fortranspecific}
The following example illustrates the use of the conditional compilation
sentinel in free form source. 
With OpenMP compilation, the conditional compilation sentinel \scode{!$} is
replaced with two spaces.

\ffreenexample{cond_comp}{1}

The following example illustrates the use of the conditional compilation 
sentinel in fixed form source. 
With OpenMP compilation, the conditional compilation sentinel \scode{!$} is
replaced with two spaces.
In fixed form source, the sentinel \scode{!$} must start on column 1 and
statements guarded by the sentinel must start after column 6.
The first line (\ucode{S-1}) is used here to identify the first 9 columns.

\fnexample{cond_comp}{1}
\end{fortranspecific}

